\chapter{L'apparato scheletrico: tronco e arti}

% ---------------------------------------------------------------------------
\section{La colonna vertebrale}

La \textbf{colonna vertebrale}, o \textbf{rachide}, è il principale sostegno
del corpo umano e dei vertebrati, ed è costituita dalla sovrapposizione di
33-34 ossa brevi. Osservata lateralmente, presenta quattro curve alternate:
la \textbf{curva cervicale}, concava posteriormente (lordosi); la
\textbf{curva toracica}, concava anteriormente (cifosi); la \textbf{curva
lombare}, concava posteriormente (lordosi); la \textbf{curva sacrale},
concava anteriormente (cifosi). Procedendo dall'alto verso il basso si
distinguono così il rachide cervicale (sede della lordosi cervicale), il
rachide dorsale (sede della cifosi dorsale), il rachide lombare (sede della
lordosi lombare) e il rachide pelvico (sede della cifosi sacrale), fino
all'osso sacro e al coccige.

\subsection{Struttura generale delle vertebre}

Ogni \textbf{vertebra} è costituita da un corpo e da un arco, che insieme
delimitano un foro. Il \textbf{corpo vertebrale} ha forma pressoché
cilindrica, con due facce, superiore e inferiore, concave nella porzione
centrale e lievemente rilevate ai bordi. Le vertebre si articolano tra loro
grazie alle sinfisi, cioè con l'interposizione di \textbf{dischi
fibrocartilaginei}, la cui funzione è rendere complementari le superfici
articolari e ammortizzare le forze di carico, statiche e dinamiche, che si
esercitano sulla colonna vertebrale; i dischi fibrocartilaginei occupano
complessivamente circa un terzo della lunghezza dell'intera colonna
vertebrale. L'\textbf{arco vertebrale} è la porzione posteriore della
vertebra, che completa il forame vertebrale; su di esso si riconoscono due
peduncoli, due masse apofisarie, due lamine e un processo spinoso.

\subsection{Vertebre cervicali (C1-C7)}

Il corpo delle vertebre cervicali aumenta di volume in direzione
craniocaudale. L'\textbf{atlante} (C1) è privo di corpo; presenta due archi,
uno anteriore e uno posteriore, e due masse apofisarie laterali, sulla cui
faccia superiore si trova la cavità glenoidea per l'articolazione con i
condili dell'osso occipitale; sull'arco anteriore si riconoscono il
tubercolo anteriore e la faccetta per il dente, su quello posteriore il
tubercolo posteriore. L'\textbf{epistrofeo}, o asse (C2), è caratterizzato
dal \textbf{dente}, che si articola con la faccetta articolare dell'arco
anteriore dell'atlante; ai lati del dente si osservano le faccette
articolari superiori. Il dente è considerato il corpo dell'atlante, rimasto
saldato a quello dell'epistrofeo. La \textbf{settima vertebra cervicale
(C7)} è una vertebra di transizione: il suo processo spinoso è lungo, non
bifido, e sporge alla base del collo, da cui il nome di \textbf{vertebra
prominente}. Nelle vertebre cervicali tipiche si riconoscono inoltre il
forame vertebrale, il forame trasverso (attraversato dall'arteria
vertebrale), il peduncolo, la lamina, il solco per il nervo spinale, il
processo uncinato e il processo trasverso.

\subsection{Vertebre toraciche (T1-T12)}

Le vertebre toraciche si articolano con le coste, presentando pertanto
faccette articolari corrispondenti. I corpi vertebrali hanno forma
pressoché cilindrica, con diametri anteroposteriore e trasverso pressoché
uguali, e aumentano gradualmente di dimensioni procedendo verso il basso.
Presentano posteriormente due faccette articolari per lato, una superiore
e una inferiore, che si articolano con le emifaccette articolari delle
teste delle coste; sul processo trasverso si trova inoltre la faccetta
costale trasversaria, per l'articolazione con il tubercolo della costa. Si
riconoscono inoltre il processo articolare superiore e inferiore, il
processo spinoso, l'arco vertebrale, il canale vertebrale e il corpo
vertebrale.

\subsection{Vertebre lombari}

Le vertebre lombari presentano un corpo vertebrale robusto, un processo
spinoso corto e quadrangolare, un ampio forame vertebrale, e processi
trasversi definiti anche processi costali; sul processo articolare
superiore si trova il processo mammillare. Le vertebre si succedono dalla
prima (L1) alla quinta (L5) aumentando progressivamente di dimensioni; tra
i corpi vertebrali si interpone lo spazio corrispondente al disco
intervertebrale, e ciascuna vertebra presenta l'incisura vertebrale
superiore e inferiore, che delimitano il forame intervertebrale, oltre
all'articolazione zigapofisaria tra i processi articolari di vertebre
contigue.

\subsection{Sacro e coccige}

L'\textbf{osso sacro} origina dalla fusione di cinque vertebre sacrali; è un
osso impari e asimmetrico, di forma piramidale quadrangolare, appiattito in
senso anteroposteriore, e si articola lateralmente con le due ossa
dell'anca, costituendo il bacino insieme al coccige. Presenta una faccia
anteriore, o pelvica, concava e liscia, una faccia posteriore, o dorsale, e
due facce laterali. Sulla faccia anteriore si trovano quattro coppie di
\textbf{forami sacrali}, disposti su due file, che comunicano con il canale
sacrale e danno passaggio ai rami anteriori dei nervi sacrali; tra i forami
si osservano le linee trasverse, residui delle saldature (sinostosi) dei
primitivi corpi vertebrali. Si riconoscono inoltre il promontorio, l'ala
del sacro, il processo articolare superiore, il canale sacrale, le creste
sacrali laterale, mediana e mediale, la superficie articolare (auricolare)
per l'osso dell'anca, la tuberosità sacrale, lo iato sacrale, le corna del
sacro e l'articolazione sacrococcigea. Il \textbf{coccige} è costituito
dalla fusione di quattro o cinque segmenti ossei, che si rimpiccioliscono
progressivamente in direzione craniocaudale.

% ---------------------------------------------------------------------------
\section{La gabbia toracica}

Le \textbf{coste} sono dodici paia, formate da una parte ossea e da una
parte cartilaginea. Le prime sette sono \textbf{coste vere}, e si
articolano con lo sterno; l'ottava, la nona e la decima sono \textbf{coste
asternali}; l'undicesima e la dodicesima sono \textbf{coste false}, poiché
non si articolano con lo sterno. L'obliquità delle coste aumenta dalla
prima alla settima e diminuisce fino alla dodicesima. Ogni costa è
costituita da un corpo e due estremità, anteriore e posteriore; il corpo è
sottile e piatto, con due facce (una esterna convessa e una interna
concava) e due margini, superiore e inferiore; l'estremità posteriore
presenta la testa, il collo e il tubercolo della costa.

Le prime due coste e le ultime due si differenziano dalle altre per alcune
caratteristiche particolari. La \textbf{prima costa} ha un raggio di
curvatura minore, è corta, spessa, appiattita, e orientata in modo da
presentare una faccia inferiore e una superiore, con margini esterno e
interno; forma il primo arco costale, diretto per tutta la sua lunghezza
verso il basso. La \textbf{seconda costa} è lunga circa il doppio della
prima e presenta un orientamento assai simile. L'\textbf{undicesima e la
dodicesima costa} si articolano, come la prima, solo con il corpo delle
vertebre corrispondenti; mancano del collo e del tubercolo e non presentano
né l'angolo costale né il solco costale.

\subsection{Sterno}

Lo \textbf{sterno} è un osso piatto, impari, in posizione mediana nella
parete anteriore del torace, che chiude anteriormente la gabbia toracica.
Si estende dall'alto al basso per circa 20 cm, corrispondendo al tratto
compreso tra T2 e T9, ed è formato da tre parti, in senso caudo-rostrale.
Il \textbf{manubrio} ha forma quadrangolare: il suo margine inferiore è
rivestito di cartilagine ialina e si articola con il margine superiore del
corpo tramite l'articolazione (sinfisi) manubrio-sternale, che consente
leggeri movimenti angolari e limitati spostamenti anteroposteriori tra
manubrio e corpo durante la respirazione. Il \textbf{corpo}, corrispondente
a T3-T8, ha l'estremità superiore articolata con il manubrio e quella
inferiore con il processo xifoideo tramite l'articolazione xifosternale
(sinfisi). Il \textbf{processo xifoideo} è la parte più piccola dello
sterno: si trova nell'epigastrio, in rapporto con il fegato, e offre
attacco ad alcuni muscoli.

% ---------------------------------------------------------------------------
\section{L'arto superiore}

L'\textbf{arto superiore} è formato da spalla, braccio, gomito, avambraccio
e mano; le sue diverse parti si articolano tra loro a livello del cingolo
scapolare, del gomito e del polso. È costituito dalla \textbf{cintura
pettorale} (scapola e clavicola, che formano la spalla) e dalla \textbf{parte
libera} (braccio, avambraccio e mano), a sua volta costituita da omero,
radio e ulna, e da carpo, metacarpo e falangi.

\subsection{Scapola e clavicola}

La \textbf{scapola} è un osso piatto, triangolare, pari e simmetrico,
situato posteriormente e superolateralmente rispetto alla gabbia toracica.
La faccia anteriore è leggermente concava, mentre quella posteriore è
divisa in due da un processo diagonale, la \textbf{spina della scapola},
che continua oltre l'angolo superolaterale con l'\textbf{acromion}.
L'angolo superolaterale presenta la \textbf{cavità glenoidea}, per
l'articolazione con l'omero.

La \textbf{clavicola} è un osso piatto e allungato, con un corpo e due
estremità, una mediale e una laterale. Il corpo ha forma a S, convesso in
avanti per i due terzi mediali e concavo in avanti nel terzo laterale; si
estende dalla base del collo fino all'apice della spalla. L'estremità
acromiale si articola con il margine mediale dell'acromion mediante
un'artrodia provvista di un disco fibrocartilagineo; una capsula molto
resistente e robusti legamenti limitano i movimenti tra le due ossa.

\subsection{Omero}

L'\textbf{omero} è un osso lungo, pari e simmetrico, che forma lo scheletro
del braccio, costituito da una diafisi (cilindrica superiormente,
triangolare inferiormente) e due epifisi. L'epifisi prossimale presenta la
testa dell'omero, il collo anatomico, il collo chirurgico, e il tubercolo
maggiore e il tubercolo minore, separati dal solco intertubercolare; dalle
creste dei due tubercoli originano rispettivamente il labbro laterale e il
labbro mediale. Il corpo presenta la tuberosità deltoidea e, in prossimità
dell'epifisi distale, si appiattisce in senso anteroposteriore,
presentando due sporgenze ossee sede di inserzioni muscolari, gli
\textbf{epicondili} laterale e mediale. Tra i due epicondili si trova una
superficie articolare irregolare, il \textbf{condilo}, caratterizzata da
due porzioni: una laterale, di forma quasi emisferica, il \textbf{capitello},
per l'articolazione con il radio, e una mediale, la \textbf{troclea}, per
l'articolazione con l'ulna; posteriormente il condilo presenta la fossa
olecranica.

Le \textbf{fratture della testa dell'omero} sono frequenti in ambito
sportivo, in particolare negli sport di contatto e in attività come sci,
snowboard, calcio e rugby.

\subsection{Radio e ulna}

L'\textbf{ulna} e il \textbf{radio} costituiscono insieme lo scheletro
dell'avambraccio; sono entrambi ossa lunghe, con due epifisi e una diafisi.
L'\textbf{ulna} è l'osso mediale, a sezione triangolare: la sua epifisi
prossimale presenta l'\textbf{incisura trocleare}, che corrisponde alla
troclea omerale ed è divisa in due settori da una cresta ossea, delimitata
dal processo coronoideo, situato in posizione antero-inferiore, e da un
voluminoso processo appuntito, l'\textbf{olecrano}, situato in posizione
postero-superiore; distalmente si trovano la testa dell'ulna e il processo
stiloideo dell'ulna. Il \textbf{radio} è laterale rispetto all'ulna, con
diafisi triangolare, incurvata verso l'ulna e a essa unita mediante la
membrana interossea. L'epifisi prossimale presenta la testa e il collo del
radio, con la tuberosità del radio; l'epifisi distale presenta lateralmente
la sporgenza del processo stiloideo del radio, medialmente l'incisura
ulnare, per l'articolazione con l'ulna, e distalmente un'ampia superficie
articolare concava, per l'articolazione con le ossa della fila prossimale
del carpo.

\subsection{La mano}

Le ossa del \textbf{carpo} si dispongono in due file. La fila prossimale
comprende, in senso lateromediale, l'osso scafoide (o navicolare), l'osso
semilunare, l'osso piramidale e il pisiforme. La fila distale comprende, in
senso lateromediale, l'osso trapezio, il trapezoide, il capitato e
l'uncinato. Il \textbf{metacarpo} è costituito da cinque ossa lunghe. Le
\textbf{falangi} costituiscono lo scheletro delle dita: sono piccole ossa
lunghe, in numero di tre (prossimale, media e distale, o ungueale) per ogni
dito, eccetto il primo, che ne possiede due; si articolano tra loro
mediante troclee, e con le ossa metacarpali tramite condili.

% ---------------------------------------------------------------------------
\section{Il bacino}

Il \textbf{bacino}, o \textbf{pelvi}, nel suo insieme è un anello osseo
costituito, in avanti e lateralmente, dalla porzione inferiore delle due
ossa dell'anca, e posteriormente dall'osso sacro e dal coccige; queste
componenti sono articolate in modo da costituire una formazione rigida che
forma la parte inferiore del tronco, con la forma di un cono irregolare. Le
pareti ossee del bacino, costituite dall'osso dell'anca, dal sacro e dal
coccige, sono rivestite da parti molli formate da tegumenti, muscoli e
fasce. La \textbf{cintura pelvica} è costituita dalle due ossa dell'anca,
articolate fra loro e unite all'osso sacro tramite l'\textbf{articolazione
sacroiliaca}. Nell'osso dell'anca si riconoscono, oltre alle strutture
descritte di seguito, il promontorio del sacro, la cresta iliaca (con
labbro interno, linea intermedia e labbro esterno), l'ala iliaca, la
grande e la piccola incisura ischiatica, la linea arcuata, la spina
ischiatica, l'eminenza ileo-pubica, il ramo superiore e inferiore del pube,
il foro otturato, il tubercolo pubico e l'arcata pubica.

\subsection{Osso dell'anca}

L'\textbf{osso dell'anca} è costituito da tre parti. L'\textbf{ileo}
comprende un corpo, lateralmente al quale si trova la superficie concava
articolare per il femore (\textbf{acetabolo}), e una parte superiore
appiattita, l'ala. L'\textbf{ischio} prende parte, come l'ileo, alla
formazione dell'acetabolo; ha un ramo a forma di "C" aperta anteriormente,
che nella parte posteriore forma la \textbf{tuberosità ischiatica}. Il
\textbf{pube}, tramite il proprio corpo fuso con quelli dell'ileo e
dell'ischio, costituisce la porzione antero-inferiore dell'acetabolo. La
\textbf{sinfisi pubica} è data dall'unione delle superfici articolari delle
porzioni pubiche delle due ossa dell'anca, fra le quali si interpone un
disco fibrocartilagineo; è un'articolazione ampia, resistente alle
tensioni elastiche.

% ---------------------------------------------------------------------------
\section{L'arto inferiore}

\subsection{Femore}

Il \textbf{femore} è l'osso dell'arto inferiore situato nella coscia, che
costituisce anche parte dell'anca e del ginocchio; è l'osso più lungo, più
voluminoso e più resistente dello scheletro. È formato da un corpo
(diafisi) e due estremità (epifisi). L'epifisi prossimale presenta una
testa cilindrica e il collo anatomico; tra il collo anatomico e la diafisi
si trovano due protuberanze per l'inserzione di muscoli, il \textbf{grande}
e il \textbf{piccolo trocantere}. La diafisi è cilindrica in sezione.
L'epifisi distale presenta gli epicondili laterale e mediale; i due
condili sono fusi anteriormente, dove presentano la faccia patellare,
mentre posteriormente sono separati da un'ampia fossa intercondilare.

La regione del femore più soggetta a \textbf{fratture} è quella del collo,
soprattutto a causa del peso corporeo a cui è sottoposta; questo problema si
acuisce in particolare in età avanzata, a causa della diminuzione della
quantità di tessuto osseo e dell'alterazione della sua composizione chimica
(\textbf{osteoporosi}). Al microscopio, il tessuto osseo osteoporotico
mostra una trama molto più rarefatta rispetto al tessuto osseo normale. La
\textbf{vitamina D} svolge un ruolo importante nel mantenimento della salute
ossea e nella prevenzione dell'osteoporosi.

\subsection{Patella}

La \textbf{patella} è il più grande osso sesamoide del corpo; è situata
nella regione anteriore del ginocchio, accolta nell'espansione tendinea del
muscolo quadricipite femorale, ed è posta anteriormente alla faccia
patellare del femore, fra i due condili. È un osso breve, di forma
triangolare, con apice inferiore e rotondeggiante, appiattito nella parte
posteriore; presenta due facce, due margini, una base e un apice; la faccia
posteriore è articolare ed è suddivisa in una porzione mediale e una
laterale.

\subsection{Tibia e perone}

La \textbf{tibia} è un osso lungo a sezione triangolare, con margine
anteriore piuttosto affilato; superiormente continua con la
\textbf{tuberosità tibiale}, su cui si inserisce il tendine del muscolo
quadricipite femorale, mentre sulla faccia posteriore è presente la linea
del muscolo soleo. L'epifisi prossimale è allargata a piramide, con la base
superiore provvista di due facce articolari (il piatto tibiale) per i
condili del femore, separate dall'eminenza intercondiloidea; distalmente si
trova il malleolo mediale. Il \textbf{perone}, o \textbf{fibula}, è un
sottile osso lungo posto lateralmente alla tibia, con la testa del perone
prossimalmente e il malleolo laterale distalmente; tibia e perone sono
uniti dalla membrana interossea e dalle articolazioni tibio-fibulari
superiore e inferiore.

\subsection{Il piede}

Le ossa del \textbf{tarso} sono sette e formano un complesso osseo che
comprende il \textbf{talo} (o astragalo), il \textbf{calcagno}, l'\textbf{osso
navicolare}, l'\textbf{osso cuboide} e le tre \textbf{ossa cuneiformi} (prima,
seconda e terza). Il talo è un osso breve coinvolto nell'articolazione
tibiotarsica, e presenta testa, collo, corpo e troclea; il calcagno
presenta il sustentaculum tali, la tuberosità calcaneale e il seno del
tarso.

Il \textbf{metatarso} è composto da cinque ossa lunghe, numerate dalla
prima alla quinta in direzione mediolaterale, incurvate verso il lato
plantare; l'estremità prossimale si articola con le ossa cuneiformi e con
l'osso cuboide per mezzo di artrodie, mentre quella distale si articola con
le falangi mediante condilartrosi. Le \textbf{falangi} del piede sono tre
per ogni dito, con l'eccezione del primo, che ne possiede due; sono piccole
ossa lunghe, articolate tra loro mediante troclee, e le falangi distali
accolgono l'unghia dorsalmente.
